% GENERATED FILE. Edit canonical lessons, then run npm run generate:books.
% canonical-source-hash: fnv1a64:755907fad1695bfa
% canonical-lessons: ZH-C351-fenxi, ZH-C351-youli, ZH-W351-li-reason, ZH-C351-jiangli, ZH-C351-lixiang, ZH-C351-xinli, ZH-R351-review-346-351

\chapter{That Makes Sense}
\label{ch:zh-that-makes-sense}
\hlchaptermodality{351}

\begin{chapteropening}
\textbf{By the end of this chapter:} \emph{I can say someone is in the right, analyse a problem, and talk about ideals and states of mind.}

Complete the last lesson of chapter 351: I can say someone is in the right, analyse a problem, and talk about ideals and states of mind.
\end{chapteropening}

\section[fēnxī]{fēnxī — to analyse}

\label{lesson:ZH-C351-fenxi}



\begin{quote}
Before the new one: say the Mandarin for to make a judgement, then the Mandarin for indeed.
\end{quote}

\subsection*{You'll want to know: fēnxī}
\emph{fēnxī} — \textquotedblleft{}to analyse\textquotedblright{}.

\begin{quote}
\emph{wǒmen yīnggāi xiān fēnxī wèntí} — We should analyse the problem first.
\end{quote}

Its characters come later; for now it is a word for the ear.

\begin{cousinweb}
\emph{fēnxī} is built from \emph{fēn} and \emph{xī}, both inherited from Middle Chinese, \emph{pjun} and \emph{sek} (Guangyun, in Baxter's transcription), which glosses the first 'to share out' and the second 'to split'.
\end{cousinweb}

\begin{grammarlens}[title={What you've built}]
One more word for giving reasons.
\end{grammarlens}

\subsection*{Your turn}
\begin{itemize}
\raggedright
  \item \emph{Say it:} \emph{fēnxī}
  \item \emph{Say it:} \emph{fēnxī}, once more
  \item \emph{Say it:} \emph{wǒmen yīnggāi xiān fēnxī wèntí} — We should analyse the problem first.
  \item \emph{Recall:} say the Mandarin for to make a judgement, then the Mandarin for indeed, then say \emph{fēnxī} again
\end{itemize}

\begin{tcolorbox}[breakable,colback=teal!4,colframe=teal!35!black,title={Before you move on}]
What does fēnxī mean? (\textquotedblleft{}To analyse\textquotedblright{}.) Say it once more.
\end{tcolorbox}
\section[\texorpdfstring{\zh{有理} (yǒulǐ)}{有理 (yǒulǐ)}]{\zh{有理} (yǒulǐ) — reasonable; in the right}

\label{lesson:ZH-C351-youli}



\begin{quote}
Before the new one: say the Mandarin for indeed, then the Mandarin for to analyse.
\end{quote}

\subsection*{You'll want to know: \zh{有理}}
\textbf{\zh{有理}} — \emph{yǒulǐ} — \textquotedblleft{}reasonable; in the right\textquotedblright{}.

\begin{quote}
\textbf{\zh{我觉得你说得有理}} \emph{wǒ juéde nǐ shuō de yǒulǐ} — I think you've got a point.
\end{quote}

The new character, in pieces: \zh{理} — eleven strokes in two pieces: on the left the jade piece, four strokes, three bars and a vertical, the last bar rising to the right; on the right \zh{里} \emph{lǐ}, seven strokes. The next lesson puts it in your hand.

Two third tones meet in this word: yǒu is said with a rising tone before lǐ, and the spelling keeps both marks.

\begin{cousinweb}
\textbf{\zh{有理}} (\emph{yǒulǐ}) is built from \zh{有} 'to have' and \zh{理} 'order', both inherited from Middle Chinese, \emph{hjuwX} and \emph{liX} (Guangyun, in Baxter's transcription), which glosses the second 'to put in order' and, quoting an older dictionary, 'to work jade'.
\end{cousinweb}

\begin{grammarlens}[title={What you've built}]
One more word for giving reasons.
\end{grammarlens}

\subsection*{Your turn}
\begin{itemize}
\raggedright
  \item \emph{Say it:} \emph{yǒulǐ}
  \item \emph{Say it:} \emph{yǒulǐ}, once more
  \item \emph{Say it:} \emph{wǒ juéde nǐ shuō de yǒulǐ} — I think you've got a point.
  \item \emph{Recall:} say the Mandarin for indeed, then the Mandarin for to analyse, then say \emph{yǒulǐ} again
\end{itemize}

\begin{tcolorbox}[breakable,colback=teal!4,colframe=teal!35!black,title={Before you move on}]
What does \zh{有理} mean? (\textquotedblleft{}Reasonable; in the right\textquotedblright{}.) Say it once more.
\end{tcolorbox}
\section[\texorpdfstring{\zh{理} (lǐ)}{理 (lǐ)}]{\zh{理} — the jade piece, then lǐ}

\label{lesson:ZH-W351-li-reason}



\begin{quote}
The word first, then the character inside it.

\begin{itemize}
\raggedright
  \item \emph{Recall:} say \emph{yǒulǐ} — \textbf{R1}, one lesson back, before it decays
\end{itemize}

\begin{quote}
\zh{里} \zh{有}
\end{quote}

\emph{Read it:} each one aloud, with nothing said first — \emph{lǐ}, \emph{yǒu}
\end{quote}

\subsection*{Script — the jade piece, then lǐ}
\hlblockfigure{\detokenize{figures/ZH-W351-li-reason-filmstrip.pdf}}{How \zh{理} is written, stroke by stroke}

\begin{quote}
\zh{理}
\end{quote}

\textbf{Two pieces, left before right.} Count the pieces, not the strokes.

\textbf{The left piece, the jade piece, four strokes}, narrow: a horizontal, a second horizontal, a vertical down through them, then a bar at the foot that rises to the right.

\textbf{The right piece, \zh{里}, seven strokes}: a vertical, a horizontal that turns down, a horizontal across the middle, a horizontal closing the box, a vertical down through the middle, a short horizontal, then the long horizontal at the foot.

4 + 7 = eleven, and the last seven your hand already knows.

\subsection*{Your turn}
\begin{itemize}
\raggedright
  \item \emph{Write it:} \zh{理}, all eleven, with the model in front of you
  \item \emph{Cover:} write it again from memory, and count the pieces, not the strokes
\end{itemize}

\begin{tcolorbox}[breakable,colback=teal!4,colframe=teal!35!black,title={Before you move on}]
How many pieces in \zh{理}? (\textbf{Two}.) What is the right piece read? (\emph{lǐ}, as in \zh{里}.) How does the jade piece end? (\textbf{With a bar that rises to the right}.)

Source: \href{https://raw.githubusercontent.com/chanind/hanzi-writer-data/68d10a4b21150cae5e1ebbd223eed289cf32d90c/data/\%E7\%90\%86.json}{Hanzi Writer Data \zh{理}, pinned stroke paths and medians}.
\end{tcolorbox}
\section[jiǎnglǐ]{jiǎnglǐ — to be reasonable; to listen to reason}

\label{lesson:ZH-C351-jiangli}



\begin{quote}
Before the new one: say the Mandarin for to analyse, then the Mandarin for in the right in an argument.
\end{quote}

\subsection*{You'll want to know: jiǎnglǐ}
\emph{jiǎnglǐ} — \textquotedblleft{}to be reasonable; to listen to reason\textquotedblright{}.

\begin{quote}
\emph{zhège rén bù jiǎnglǐ} — This man won't listen to reason.
\end{quote}

Its characters come later; for now it is a word for the ear.

Two third tones meet in this word: jiǎng is said with a rising tone before lǐ, and the spelling keeps both marks.

\begin{cousinweb}
\emph{jiǎnglǐ} is built from \emph{jiǎng} and \zh{理} 'order', both inherited from Middle Chinese, \emph{kaewngX} and \emph{liX} (Guangyun, in Baxter's transcription), which glosses the first 'to tell' and 'to discuss', and the second 'to put in order'.
\end{cousinweb}

\begin{grammarlens}[title={What you've built}]
One more word for giving reasons.
\end{grammarlens}

\subsection*{Your turn}
\begin{itemize}
\raggedright
  \item \emph{Say it:} \emph{jiǎnglǐ}
  \item \emph{Say it:} \emph{jiǎnglǐ}, once more
  \item \emph{Say it:} \emph{zhège rén bù jiǎnglǐ} — This man won't listen to reason.
  \item \emph{Recall:} say the Mandarin for to analyse, then the Mandarin for in the right in an argument, then say \emph{jiǎnglǐ} again
\end{itemize}

\begin{tcolorbox}[breakable,colback=teal!4,colframe=teal!35!black,title={Before you move on}]
What does jiǎnglǐ mean? (\textquotedblleft{}To be reasonable; to listen to reason\textquotedblright{}.) Say it once more.
\end{tcolorbox}
\section[\texorpdfstring{\zh{理想} (lǐxiǎng)}{理想 (lǐxiǎng)}]{\zh{理想} (lǐxiǎng) — ideal; aspiration}

\label{lesson:ZH-C351-lixiang}



\begin{quote}
Before the new one: say the Mandarin for in the right in an argument, then the Mandarin for to listen to reason.
\end{quote}

\subsection*{You'll want to know: \zh{理想}}
\textbf{\zh{理想}} — \emph{lǐxiǎng} — \textquotedblleft{}ideal; aspiration\textquotedblright{}.

\begin{quote}
\textbf{\zh{我的理想是当医生}} \emph{wǒ de lǐxiǎng shì dāng yīshēng} — My dream is to be a doctor.
\end{quote}

Two third tones meet in this word: lǐ is said with a rising tone before xiǎng, and the spelling keeps both marks.

\begin{cousinweb}
\textbf{\zh{理想}} (\emph{lǐxiǎng}) is built from \zh{理} 'order' and \zh{想} 'to think', both inherited from Middle Chinese, \emph{liX} and \emph{sjangX} (Guangyun, in Baxter's transcription), which glosses the first 'to put in order' and the second 'to think'.
\end{cousinweb}

\begin{grammarlens}[title={What you've built}]
One more word for giving reasons.
\end{grammarlens}

\subsection*{Your turn}
\begin{itemize}
\raggedright
  \item \emph{Say it:} \emph{lǐxiǎng}
  \item \emph{Say it:} \emph{lǐxiǎng}, once more
  \item \emph{Say it:} \emph{wǒ de lǐxiǎng shì dāng yīshēng} — My dream is to be a doctor.
  \item \emph{Recall:} say the Mandarin for in the right in an argument, then the Mandarin for to listen to reason, then say \emph{lǐxiǎng} again
\end{itemize}

\begin{tcolorbox}[breakable,colback=teal!4,colframe=teal!35!black,title={Before you move on}]
What does \zh{理想} mean? (\textquotedblleft{}Ideal; aspiration\textquotedblright{}.) Say it once more.
\end{tcolorbox}
\section[\texorpdfstring{\zh{心理} (xīnlǐ)}{心理 (xīnlǐ)}]{\zh{心理} (xīnlǐ) — state of mind; psychology}

\label{lesson:ZH-C351-xinli}



\begin{quote}
Before the new one: say the Mandarin for to listen to reason, then the Mandarin for an ideal to live up to.
\end{quote}

\subsection*{You'll want to know: \zh{心理}}
\textbf{\zh{心理}} — \emph{xīnlǐ} — \textquotedblleft{}state of mind; psychology\textquotedblright{}.

\begin{quote}
\emph{wǒ juéde zhè shì xīnlǐ wèntí} — I think it's all in the mind.
\end{quote}

\begin{cousinweb}
\textbf{\zh{心理}} (\emph{xīnlǐ}) is built from \zh{心} 'the heart' and \zh{理} 'order', both inherited from Middle Chinese, \emph{sim} and \emph{liX} (Guangyun, in Baxter's transcription), which glosses the second 'to put in order'.
\end{cousinweb}

\begin{grammarlens}[title={What you've built}]
One more word for giving reasons.
\end{grammarlens}

\subsection*{Your turn}
\begin{itemize}
\raggedright
  \item \emph{Say it:} \emph{xīnlǐ}
  \item \emph{Say it:} \emph{xīnlǐ}, once more
  \item \emph{Say it:} \emph{wǒ juéde zhè shì xīnlǐ wèntí} — I think it's all in the mind.
  \item \emph{Recall:} say the Mandarin for to listen to reason, then the Mandarin for an ideal to live up to, then say \emph{xīnlǐ} again
\end{itemize}

\begin{tcolorbox}[breakable,colback=teal!4,colframe=teal!35!black,title={Before you move on}]
What does \zh{心理} mean? (\textquotedblleft{}State of mind; psychology\textquotedblright{}.) Say it once more.
\end{tcolorbox}
\section[(dialogue)]{Looking back: from Thinking It Over to That Makes Sense}

\label{lesson:ZH-R351-review-346-351}



\begin{quote}
Say the Mandarin for to think something over and a point of view.
\end{quote}

\subsection*{Your turn}
\begin{itemize}
\raggedright
  \item \emph{Say it:} to think something over, a point of view, your impressions afterwards, to wrongly assume, to make a rough guess
  \item \emph{Say it:} a bright idea, what something means, to think hard, the significance of something, to know from experience
  \item \emph{Say it:} to discuss, to talk it over together, to argue over a point, the other party, all of one mind
  \item \emph{Say it:} to regard, a way of thinking, one version of events, the way something is done, to argue in favour of
  \item \emph{Say it:} a fact, to prove, a basis for a claim, to make a judgement, indeed
  \item \emph{Say it:} to analyse, in the right in an argument, to listen to reason, an ideal to live up to, a person's state of mind
  \item \emph{Read it:} \zh{思} \zh{法} \zh{理} — each one aloud, then a word that holds it
\end{itemize}

\begin{tcolorbox}[breakable,colback=teal!4,colframe=teal!35!black,title={Before you move on}]
To think something over? (\emph{kǎolǜ}.) A point of view? (\emph{guāndiǎn}.) Your impressions afterwards? (\textbf{\zh{感想}}, \emph{gǎnxiǎng}.) To wrongly assume? (\textbf{\zh{以为}}, \emph{yǐwéi}.) To make a rough guess? (\emph{gūjì}.) A bright idea? (\emph{zhǔyi}.) What something means? (\textbf{\zh{意思}}, \emph{yìsi}.) To think hard? (\emph{sīkǎo}.) The significance of something? (\emph{yìyì}.) To know from experience? (\textbf{\zh{体会}}, \emph{tǐhuì}.) To discuss? (\emph{tǎolùn}.) To talk it over together? (\emph{shāngliang}.) To argue over a point? (\emph{zhēnglùn}.) The other party? (\emph{duìfāng}.) All of one mind? (\emph{yízhì}.) To regard? (\emph{kàndài}.) A way of thinking? (\textbf{\zh{想法}}, \emph{xiǎngfǎ}.) One version of events? (\textbf{\zh{说法}}, \emph{shuōfa}.) The way something is done? (\textbf{\zh{做法}}, \emph{zuòfǎ}.) To argue in favour of? (\emph{zhǔzhāng}.) A fact? (\emph{shìshí}.) To prove? (\emph{zhèngmíng}.) A basis for a claim? (\emph{gēnjù}.) To make a judgement? (\emph{pànduàn}.) Indeed? (\emph{quèshí}.) To analyse? (\emph{fēnxī}.) In the right in an argument? (\textbf{\zh{有理}}, \emph{yǒulǐ}.) To listen to reason? (\emph{jiǎnglǐ}.) An ideal to live up to? (\textbf{\zh{理想}}, \emph{lǐxiǎng}.) A person's state of mind? (\textbf{\zh{心理}}, \emph{xīnlǐ}.)
\end{tcolorbox}
